\bibitem{K17Carayol1986}
H.~Carayol, \emph{Sur les repr\'esentations $\ell$-adiques associ\'ees
aux formes modulaires de Hilbert}, Annales scientifiques de l'\'Ecole
Normale Sup\'erieure (4) \textbf{19} (1986), 409--468.
\S0.3, p.~409, and \S0.7, Theorem~A, pp.~410--411.
\url{https://www.numdam.org/item/10.24033/asens.1512.pdf}.

\bibitem{K17SaitoSurvey}
T.~Saito, \emph{Galois representations in arithmetic geometry},
author-hosted survey, \S2b--2c, pp.~7--9,
especially Theorem~3, p.~9.
\url{https://www.ms.u-tokyo.ac.jp/~t-saito/pp/suams.pdf}.
Theorem~3 explicitly attributes the classical modular-form compatibility
theorem to reference [S7], T.~Saito,
\emph{Modular forms and $p$-adic Hodge theory},
Inventiones Mathematicae \textbf{129} (1997), 607--620.
The exact theorem in this author survey was read and visually verified;
the full 1997 publisher PDF was not recovered in this run.

\bibitem{K17ColmezFontaine2000}
P.~Colmez and J.-M.~Fontaine,
\emph{Construction des repr\'esentations $p$-adiques semi-stables},
Inventiones Mathematicae \textbf{140} (2000), 1--43.
Author version: \S3.3, p.~21; \S4.1, Proposition~4.2, p.~23;
\S4.2, Proposition~4.4, p.~24.
\url{https://webusers.imj-prg.fr/~pierre.colmez/CF.pdf}.
